\documentclass[11pt,a4paper]{article}

\usepackage[utf8]{inputenc}
\usepackage[T1]{fontenc}
\usepackage[french]{babel}
\usepackage{lmodern}
\usepackage[a4paper,margin=2.5cm]{geometry}
\usepackage{hyperref}
\usepackage{xcolor}
\usepackage{listings}
\usepackage{array}
\usepackage{tabularx}
\usepackage{longtable}
\usepackage{booktabs}
\usepackage{enumitem}
\usepackage{amsmath,amssymb}
\usepackage{tcolorbox}
\usepackage{fancyvrb}
\usepackage{titlesec}

\definecolor{accent}{HTML}{2A5A8A}
\definecolor{muted}{HTML}{555555}
\definecolor{bgcode}{HTML}{F4F4F4}
\definecolor{warn}{HTML}{B85C00}
\definecolor{ok}{HTML}{2A7A2A}
\definecolor{bordercol}{HTML}{DDDDDD}
\definecolor{warnbg}{HTML}{FFF8E6}
\definecolor{okbg}{HTML}{EEF9EE}
\definecolor{thbg}{HTML}{EEF3F8}

\hypersetup{
  colorlinks=true,
  linkcolor=accent,
  urlcolor=accent,
  citecolor=accent
}

\titleformat{\section}
  {\color{accent}\Large\bfseries}
  {\thesection}{1em}{}
  [\vspace{-0.4em}\color{bordercol}\titlerule]
\titleformat{\subsection}
  {\color{accent}\large\bfseries}
  {\thesubsection}{1em}{}

\lstdefinestyle{codestyle}{
  backgroundcolor=\color{bgcode},
  basicstyle=\ttfamily\small,
  breaklines=true,
  frame=leftline,
  framerule=1pt,
  rulecolor=\color{accent},
  xleftmargin=1em,
  framexleftmargin=1em,
  showstringspaces=false,
  columns=fullflexible,
  keepspaces=true
}
\lstset{style=codestyle}

\newtcolorbox{notebox}{
  colback=warnbg,
  colframe=warn,
  boxrule=0pt,
  leftrule=4pt,
  arc=2pt,
  left=8pt,right=8pt,top=6pt,bottom=6pt
}
\newtcolorbox{okbox}{
  colback=okbg,
  colframe=ok,
  boxrule=0pt,
  leftrule=4pt,
  arc=2pt,
  left=8pt,right=8pt,top=6pt,bottom=6pt
}
\newtcolorbox{formulabox}{
  colback=white,
  colframe=bordercol,
  boxrule=1pt,
  arc=2pt,
  left=8pt,right=8pt,top=6pt,bottom=6pt,
  halign=center
}

\title{\color{accent}\textbf{Calcul de la vérité terrain pour RTS-LLM}}
\author{\textcolor{muted}{\small Document méthodologique --- projet de mémoire sur la sélection de tests de régression assistée par LLM}}
\date{}

\begin{document}

\maketitle
\thispagestyle{empty}

\renewcommand{\contentsname}{Sommaire}
\tableofcontents
\vspace{1em}
\hrule
\vspace{1em}

\section{Pourquoi une vérité terrain ?}
\label{sec:intro}

RTS-LLM est un outil de \textbf{sélection de tests de régression} : à partir d'un diff Git, il \emph{prédit} quels scénarios BDD doivent être ré-exécutés. Pour évaluer cette prédiction, il faut un \textbf{oracle} qui dit, pour le même diff, quels scénarios \emph{auraient réellement dû} être ré-exécutés.

Cet oracle s'appelle la \textbf{vérité terrain}. Sa fonction n'est pas d'être un outil RTS concurrent, mais un étalon de mesure :

\begin{Verbatim}[frame=single,fontsize=\small]
                  Diff Git d'un commit
                          |
              +-----------+----------+
              |                      |
              v                      v
    RTS-LLM (statique + LLM)   Vérité terrain (JaCoCo)
    -> S = {sc_1, sc_5}         -> G = {sc_1, sc_3}
              |                      |
              +----------+-----------+
                         |
                         v
                  Métriques :
                   Précision = |S inter G| / |S|
                   Rappel    = |S inter G| / |G|
\end{Verbatim}

La vérité terrain est \textbf{par définition correcte au regard de l'exécution réelle des tests} ; les écarts mesurés (faux positifs, faux négatifs de RTS-LLM) reflètent les limites de la prédiction, pas de l'oracle.


\section{Principe : couverture dynamique avec JaCoCo}
\label{sec:principe}

La méthodologie repose sur l'observation suivante : \emph{un scénario BDD est impacté par une modification si, et seulement si, son exécution traverse au moins une des méthodes modifiées}.

La \textbf{couverture dynamique} répond exactement à cette question : pour chaque scénario, on enregistre les méthodes effectivement exécutées lors de son passage, en instrumentant le bytecode avec \href{https://www.jacoco.org/}{JaCoCo}.

L'équation centrale de la vérité terrain s'écrit alors :

\begin{formulabox}
$G(\mathit{diff}) = \{\, s \in \text{Scénarios} \mid \text{Couverture}(s) \cap \text{MéthodesModifiées}(\mathit{diff}) \neq \emptyset \,\}$
\end{formulabox}

\noindent où :
\begin{itemize}
  \item \textbf{Couverture(s)} est l'ensemble des méthodes exécutées par le scénario \emph{s}, mesuré une fois pour toutes par JaCoCo dans la version de référence du code ;
  \item \textbf{MéthodesModifiées(diff)} est l'ensemble des méthodes que le diff Git modifie, calculé par croisement entre les \emph{hunks} du diff et les plages de méthodes issues de l'AST.
\end{itemize}


\section{Pipeline de calcul, étape par étape}
\label{sec:pipeline}

\subsection{Préparation du projet cible}

Le projet à évaluer doit avoir le plugin \texttt{jacoco-maven-plugin} configuré, ou bien être lancé avec l'agent JaCoCo injecté à la ligne de commande (mode \texttt{--no-pom-changes}, qui évite de polluer l'historique Git).

\begin{lstlisting}[language=XML]
<!-- Dans le pom.xml du projet cible (mode permanent) -->
<plugin>
  <groupId>org.jacoco</groupId>
  <artifactId>jacoco-maven-plugin</artifactId>
  <version>0.8.11</version>
  <executions>
    <execution><goals><goal>prepare-agent</goal></goals></execution>
  </executions>
</plugin>
\end{lstlisting}

\subsection{Étape 1 --- Collecte de la couverture par scénario}

Le script \texttt{scripts/run\_per\_scenario\_coverage.sh} exécute chaque scénario Cucumber dans une JVM distincte, puis renomme le fichier \texttt{target/jacoco.exec} produit en \texttt{sc\_NNN.exec}.

\begin{lstlisting}[language=bash]
./scripts/run_per_scenario_coverage.sh --no-pom-changes /chemin/projet
\end{lstlisting}

Pour chaque scénario, la commande Maven exécutée est de la forme :

\begin{lstlisting}[language=bash]
mvn test \
  -Dcucumber.filter.name="^Nom du scenario$" \
  -DargLine="-javaagent:jacocoagent.jar=destfile=target/jacoco.exec,append=false"
\end{lstlisting}

Sortie produite :

\begin{lstlisting}
coverage-per-scenario/
|-- sc_001.exec        // couverture binaire JaCoCo
|-- sc_002.exec
|-- ...
`-- index.txt          // sc_001|Should be authenticated
                       // sc_002|Add bruce wayne user
                       // ...
\end{lstlisting}

\begin{notebox}
\textbf{Choix d'isolation} --- un fork JVM par scénario élimine la pollution d'état (caches statiques, sessions persistantes, base in-memory). Le coût est l'overhead de démarrage Maven, négligeable à l'échelle d'une évaluation académique.
\end{notebox}

\subsection{Étape 2 --- Construction de la map scénario \(\rightarrow\) méthodes}

La classe \texttt{GroundTruthBuilder} lit chaque \texttt{.exec} via l'API JaCoCo (\texttt{ExecFileLoader} + \texttt{Analyzer}), en croisant les probes activées avec le bytecode compilé. Pour chaque scénario, elle produit la liste des méthodes effectivement exécutées.

\begin{lstlisting}[language=bash]
java -cp target/rts-llm-1.0-SNAPSHOT.jar com.rts.llm.GroundTruthBuilder build \
    /chemin/projet/coverage-per-scenario \
    /chemin/projet/target/classes \
    /chemin/projet/target/test-classes \
    coverage.json
\end{lstlisting}

Format de l'identifiant méthode : \texttt{FQN.NomClasse.nomMéthode}, identique à celui produit par \texttt{StaticAnalyzer}, pour permettre l'intersection avec les méthodes modifiées extraites du diff.

Normalisations appliquées sur la sortie JaCoCo :
\begin{itemize}
  \item \texttt{com/example/Foo} \(\rightarrow\) \texttt{com.example.Foo} (slash \(\rightarrow\) point)
  \item \texttt{Outer\$Inner} \(\rightarrow\) \texttt{Outer.Inner} (classes imbriquées)
  \item \texttt{lambda\$importFrom\$0} \(\rightarrow\) \texttt{importFrom} (repli des lambdas sur leur méthode englobante)
  \item \texttt{<init>}, \texttt{<clinit>}, \texttt{access\$N} : ignorés (non trackés par \texttt{StaticAnalyzer})
\end{itemize}

Extrait de \texttt{coverage.json} :

\begin{lstlisting}[language=Java]
{
  "sc_001": {
    "name": "Should be authenticated",
    "methods": [
      "fr.redfroggy.bdd.restapi.user.UserController.getAll",
      "fr.redfroggy.bdd.restapi.glue.AbstractBddStepDefinition.setUp",
      ...
    ]
  },
  ...
}
\end{lstlisting}

\subsection{Étape 3 --- Extraction des méthodes modifiées par le diff}

Pour un diff Git donné (\texttt{git diff C\^{} C}), \texttt{SymbolExtractor} :

\begin{enumerate}
  \item parse les en-têtes de hunks (\texttt{@@ -X +Y,Z @@}) pour récupérer les plages de lignes modifiées dans chaque fichier ;
  \item croise ces plages avec \texttt{methodRangesByFile} (calculé par \texttt{StaticAnalyzer} via JavaParser) ;
  \item retourne la liste des méthodes dont la plage chevauche au moins un hunk.
\end{enumerate}

\subsection{Étape 4 --- Intersection finale}

La vérité terrain pour le diff est obtenue par intersection :

\begin{lstlisting}[language=Java]
for (Map.Entry<String, ScenarioCoverage> e : coverageMap.entrySet()) {
    if (!Collections.disjoint(e.getValue().methods, modifiedMethods)) {
        groundTruth.add(e.getKey());
    }
}
\end{lstlisting}

\subsection{Boucle sur l'historique Git}

Le script \texttt{scripts/eval\_history.sh} automatise les étapes 1 à 4 sur une plage de commits :

\begin{Verbatim}[frame=single,fontsize=\small]
Pour chaque commit C dans la plage :

  1. checkout C^ (parent)         <- état AVANT le diff
  2. mvn test-compile             <- bytecode pour analyse JaCoCo
  3. run_per_scenario_coverage.sh <- un .exec par scenario
  4. GroundTruthBuilder build     <- coverage.json
  5. checkout C
  6. GroundTruthBuilder eval      <- calcule G pour le diff C^ -> C
  7. agrege dans un CSV
\end{Verbatim}

\begin{notebox}
\textbf{Important} --- la couverture est collectée sur le \emph{parent} du commit (\texttt{C\^{}}), c'est-à-dire l'état du code \emph{avant} les modifications. C'est nécessaire pour que les noms et plages des méthodes correspondent à celles que le diff référence en \texttt{-}.
\end{notebox}


\section{Métriques d'évaluation}
\label{sec:metrics}

Pour un commit donné, on compare deux ensembles :
\begin{itemize}
  \item \textbf{G} --- la vérité terrain (scénarios réellement impactés)
  \item \textbf{S} --- la sélection produite par RTS-LLM
\end{itemize}

\begin{center}
\begin{tabularx}{\textwidth}{|l|l|l|X|}
\hline
\rowcolor{thbg}
\textbf{Sigle} & \textbf{Nom} & \textbf{Définition} & \textbf{Interprétation} \\
\hline
\textbf{TP} & True Positive  & $|S \cap G|$         & Sélectionnés et impactés --- bonne décision \\
\hline
\textbf{FP} & False Positive & $|S \setminus G|$    & Sélectionnés mais pas impactés --- temps perdu \\
\hline
\textbf{FN} & False Negative & $|G \setminus S|$    & Impactés mais oubliés --- \textcolor{warn}{danger} \\
\hline
\end{tabularx}
\end{center}

\begin{formulabox}
$\text{Précision} = \dfrac{TP}{TP + FP} \qquad
 \text{Rappel} = \dfrac{TP}{TP + FN} \qquad
 F_1 = \dfrac{2 \cdot P \cdot R}{P + R}$
\end{formulabox}

Pour un outil RTS, le \textbf{rappel est l'indicateur prioritaire} : un FN signifie qu'un test impacté n'est pas exécuté, donc qu'un bug peut atteindre la production sans alerte. Un FP n'est qu'un test exécuté inutilement.


\section{Limitations connues}
\label{sec:limites}

\subsection{Surcharges de méthodes (overloads)}

Les identifiants utilisés (\texttt{FQN.Class.method}) \textbf{n'incluent pas la signature}. Si une classe possède plusieurs surcharges et que seule l'une d'entre elles est modifiée, la vérité terrain marquera comme TP des scénarios qui n'utilisaient en réalité qu'une autre surcharge.

\textbf{Effet} : surestime \texttt{gt\_size}. Faible impact sur les projets sans surcharges intensives.

\subsection{Méthodes ajoutées par le diff}

La couverture est mesurée à \texttt{C\^{}}. Si \texttt{C} introduit une méthode neuve, aucun scénario n'a pu la \og couvrir \fg{} dans la map produite à \texttt{C\^{}}. La vérité terrain ignore donc les scénarios qui, après \texttt{C}, utiliseront cette méthode.

\textbf{Effet} : sous-estime \texttt{gt\_size}, biaise contre RTS-LLM s'il a (correctement) sélectionné ces scénarios.

\emph{Atténuation possible} : doubler le coût de collecte en mesurant aussi la couverture à \texttt{C} et en unissant les deux maps.

\subsection{Classes anonymes}

JaCoCo nomme une méthode dans une classe anonyme \texttt{EnclosingClass\$1.method}. Après normalisation : \texttt{EnclosingClass.1.method}. JavaParser, lui, ne génère pas de \texttt{ClassOrInterfaceDeclaration} pour les classes anonymes : \texttt{StaticAnalyzer} ne produit donc aucun identifiant correspondant. \textbf{Méthode invisible des deux côtés}.

\subsection{Tests qui échouent}

Le script garde le fichier \texttt{.exec} même si \texttt{mvn test} retourne un code d'erreur. Si un scénario plante avant d'atteindre le code modifié (setup cassé, dépendance externe indisponible), sa couverture est incomplète \(\rightarrow\) sous-estime \texttt{gt\_size}.

\subsection{Renommages}

Un renommage de méthode dans \texttt{C} apparaît comme une suppression dans \texttt{C\^{}} et un ajout dans \texttt{C}. \texttt{SymbolExtractor} peut ne pas corréler les deux. Heureusement, un renommage pur n'altère pas le comportement observable : aucun scénario ne devrait être impacté de toute façon.


\section{Sanity checks de validation}
\label{sec:sanity}

Avant de se fier aux chiffres, quatre vérifications rapides :

\begin{center}
\begin{tabularx}{\textwidth}{|l|X|X|}
\hline
\rowcolor{thbg}
\textbf{Check} & \textbf{Procédure} & \textbf{Attendu} \\
\hline
Diff trivial & Commit qui ne touche que des commentaires & \texttt{gt\_size = 0} \\
\hline
Diff massif & Commit qui modifie une classe centrale & \texttt{gt\_size} proche du total \\
\hline
Présence dans la map & Vérifier qu'au moins une méthode du diff figure dans \texttt{coverage.json} & \texttt{grep "NomMéthode" coverage.json} retourne quelque chose \\
\hline
Cross-check avant/après & Sur 2-3 commits : \texttt{mvn test} à \texttt{C\^{}} puis à \texttt{C}, comparer les scénarios qui changent de résultat & L'ensemble obtenu doit être \textbf{inclus} dans \texttt{G} \\
\hline
\end{tabularx}
\end{center}

\begin{okbox}
\textbf{Recommandation} --- exécuter le quatrième check sur un échantillon est l'argument le plus solide à donner dans un mémoire pour défendre la fiabilité de la vérité terrain.
\end{okbox}


\section{Bibliographie \& rapport au mémoire}
\label{sec:bibliographie}

La méthodologie ici décrite est une instance du paradigme \emph{safe RTS} par couverture dynamique, formalisé notamment par :

\begin{itemize}
  \item Rothermel \& Harrold (1997), \emph{A Safe, Efficient Regression Test Selection Technique}, ACM TOSEM.
  \item Gligoric, Eloussi \& Marinov (2015), \emph{Practical Regression Test Selection with Dynamic File Dependencies} --- outil \textbf{Ekstazi}.
  \item Legunsen et al. (2017), \emph{An Extensive Study of Static Regression Test Selection in Modern Software Evolution} --- outil \textbf{STARTS}.
\end{itemize}

Ces travaux utilisent JaCoCo (ou des instrumentations équivalentes) comme oracle, exactement de la manière décrite dans ce document. Ils servent de référence à la fois \emph{baseline} (rappel = 1 par construction quand l'instrumentation est complète) et \emph{justification méthodologique} de l'approche.

Pour une évaluation de RTS-LLM, la vérité terrain est l'objet de comparaison ; les métriques de précision, rappel et F1 publiées dans le mémoire sont la projection directe des intersections $S \cap G$, $S \setminus G$, $G \setminus S$ calculées par \texttt{GroundTruthBuilder.evaluate()}.

\vspace{2em}
\hrule
\vspace{0.5em}
\begin{center}
\textcolor{muted}{\small Généré pour le projet RTS-LLM --- \texttt{scripts/run\_per\_scenario\_coverage.sh}, \texttt{scripts/eval\_history.sh}, \texttt{com.rts.llm.GroundTruthBuilder}.}
\end{center}

\end{document}
